\chapter{数列}
\setcounter{choicecounter}{0}

%% ---- 知识框架 ---- %%
\begin{knowledgeframe}
\begin{itemize}
	\item \textbf{数列的通项与前 $n$ 项和}：$S_n = a_1 + a_2 + \cdots + a_n$；$a_n = S_n - S_{n-1}$（$n \ge 2$）。
	\item \textbf{等差数列}：$a_n = a_1 + (n-1)d$，$S_n = \frac{n(a_1+a_n)}{2}$。
	\item \textbf{等比数列}：$a_n = a_1 q^{n-1}$，$S_n = \frac{a_1(1-q^n)}{1-q}$（$q \ne 1$）。
	\item \textbf{递推关系的处理}：构造辅助数列（如取倒数）将递推式转化为等比数列。
	\item \textbf{单调性判断}：$a_{n+1}-a_n$ 的符号或 $\frac{a_{n+1}}{a_n}$ 与 1 的大小比较。
\end{itemize}
\end{knowledgeframe}

%% ---- 第一节：数列的性质与运算 ---- %%
\section{数列的性质与运算}

\startexercise

\subsection*{A组}

\begin{exercise}[选修二P9习题]
	已知函数 $f(x) = \dfrac{2^x - 1}{2^x}$（$x \in \mathbb{R}$），设数列 $\{a_n\}$ 的通项公式为 $a_n = f(n)$（$n \in \mathbb{N}^*$）。
	\begin{enumerate}
		\item[(1)] 求证：$a_n \ge \dfrac{1}{2}$；
		\item[(2)] $\{a_n\}$ 是递增数列还是递减数列？为什么？
	\end{enumerate}
\end{exercise}

\subsection*{B组}

\begin{exercise}[选修二P41拓广探索]
	已知数列 $\{a_n\}$ 的首项 $a_1 = \dfrac{3}{5}$，且满足 $a_{n+1} = \dfrac{3a_n}{2a_n+1}$。
	\begin{enumerate}
		\item[(1)] 求证：数列 $\left\{\dfrac{1}{a_n} - 1\right\}$ 为等比数列；
		\item[(2)] 若 $\dfrac{1}{a_1} + \dfrac{1}{a_2} + \dfrac{1}{a_3} + \cdots + \dfrac{1}{a_n} < 100$，求满足条件的最大整数 $n$。
	\end{enumerate}
\end{exercise}

\subsection*{C组}

\begin{exercise}[选修二P41页]
	已知数列 $\{a_n\}$ 为等差数列，$a_1 = 1$，$a_3 = 2\sqrt{2}+1$，前 $n$ 项和为 $S_n$，数列 $\{b_n\}$ 满足 $b_n = \dfrac{S_n}{n}$。求证：
	\begin{enumerate}
		\item[(1)] 数列 $\{b_n\}$ 为等差数列；
		\item[(2)] 数列 $\{a_n\}$ 中的任意三项均不能构成等比数列。
	\end{enumerate}
\end{exercise}

\stopexercise
